\documentclass[a4paper,12pt]{article}
\usepackage[T2A]{fontenc}
\usepackage[utf8]{inputenc}
\usepackage[russian]{babel}
\usepackage{amsmath,amssymb,mathtools}
\usepackage{helvet}
\renewcommand{\familydefault}{\sfdefault}
\usepackage{icomma}
\usepackage{geometry}
\geometry{a4paper,left=22mm,right=22mm,top=22mm,bottom=21mm,headheight=26pt}
\usepackage{microtype}
\usepackage{tikz}
\usetikzlibrary{calc,arrows.meta,positioning,shapes.geometric}
\usepackage{tabularx,booktabs}
\usepackage{needspace}
\usepackage{enumitem}
\usepackage{hyperref}
\usepackage{parskip}
\usepackage{fancyhdr}
\usepackage{setspace}
\usepackage{xcolor}
\definecolor{accent}{HTML}{147777}
\definecolor{darkaccent}{HTML}{184E54}
\definecolor{textgray}{HTML}{41464A}
\definecolor{rulegray}{HTML}{CAD6D7}
\hypersetup{colorlinks=true,urlcolor=darkaccent,linkcolor=darkaccent,pdfauthor={\texorpdfstring{Артём Александрович Лёвин}{Artem Levin}},pdftitle={Вклады. Часть 2 — чек-лист, 15.10.26}}
\onehalfspacing
\setlength{\parindent}{0pt}
\setlist[itemize]{leftmargin=6mm,itemsep=1pt,topsep=3pt}
\setlist[enumerate]{leftmargin=8mm,itemsep=7pt,topsep=4pt}
\pagestyle{fancy}
\fancyhf{}
\fancyhead[L]{\small\color{darkaccent}Вклады. Часть 2}
\fancyhead[R]{\small\color{textgray}15.10.2026}
\fancyfoot[C]{\small\color{textgray}\thepage}
\renewcommand{\headrulewidth}{0.3pt}
\renewcommand{\headrule}{\hbox to\headwidth{\color{rulegray}\leaders\hrule height \headrulewidth\hfill}}
\newcommand{\major}[1]{\Needspace{7\baselineskip}\vspace{5mm}\noindent{\large\bfseries\color{darkaccent}#1}\par\vspace{1mm}}
\newcommand{\minor}[1]{\vspace{2mm}\noindent{\bfseries\color{accent}#1}\par}
\newcommand{\checkitem}[1]{\noindent$\square$\enspace #1\par}
\newcommand{\mone}[1]{\text{#1}}
\begin{document}
\thispagestyle{empty}
\vspace*{31mm}
\begin{center}
{\color{accent}\large МАТЕМАТИКА \quad / \quad ФИНАНСОВЫЕ ЗАДАЧИ}\par
\vspace{16mm}
{\fontsize{32}{37}\selectfont\bfseries Чек-лист}\par
\vspace{6mm}
{\fontsize{25}{31}\selectfont\bfseries Вклады. Часть 2}\par
\vspace{10mm}
{\large Сложные проценты, пополнения, снятия\\[2mm]и переменная процентная ставка}\par
\vspace{24mm}
{\large 15 октября 2026 года}\par
\vfill
{\color{darkaccent}\normalsize Лёвин Артём Александрович}\par
{\color{textgray}\normalsize эксклюзивно для Матвея Горбачева}\par
\vspace{16mm}
\end{center}
\begin{tikzpicture}[remember picture,overlay]
\draw[accent!54,line width=1.2pt] ($(current page.south west)+(24mm,22mm)$)
 .. controls ($(current page.south west)+(80mm,33mm)$) and ($(current page.south east)+(-78mm,9mm)$) .. ($(current page.south east)+(-24mm,22mm)$);
\end{tikzpicture}
\clearpage

\major{1. Что нужно уметь после занятия}
\checkitem{Отличать первоначальную сумму, процентную ставку и срок вклада.}
\checkitem{Рассчитывать капитализацию по годам и применять формулу сложных процентов.}
\checkitem{Учитывать пополнения и снятия \emph{до} очередного начисления процентов.}
\checkitem{Составлять таблицу и буквенную модель для трёхлетнего вклада.}
\checkitem{Находить неизвестные начальную сумму, ставку или срок.}
\checkitem{Выбирать между готовой формулой и расчётом по годам.}

\major{2. Обозначения и основной принцип}
Пусть $S>0$ --- первоначальная сумма вклада; $P$ --- годовая ставка \emph{в процентах};
$n$ --- число лет. Введём долю ставки $p$ и повышающий коэффициент $k$:
\[
 p=\frac{P}{100},\qquad k=1+p=1+\frac{P}{100}.
\]
Например, при ставке $15\%$ (то есть $P=15$) получаем $p=0{,}15$ и $k=1{,}15$.

\minor{Капитализация}
В конце года банк начисляет проценты на сумму, находящуюся на вкладе \emph{после операций в начале этого года}. Поэтому сумма в конце года становится начальной суммой следующего года \emph{до нового пополнения или снятия}.

Если никаких дополнительных операций нет и ставка неизменна:
\[
\boxed{\quad S_n=Sk^n=S\left(1+\frac{P}{100}\right)^n.\quad}
\]
Это \textbf{формула сложных процентов}: проценты каждого года в дальнейшем сами приносят проценты.

\minor{Мини-пример: вклад без дополнительных операций}
$S=100$ тыс. руб., ставка $10\%$, $n=3$:
\[
100\longrightarrow110\longrightarrow121\longrightarrow133{,}1
\quad (\mone{тыс. руб.}).
\]
Проверка по формуле: $S_3=100\cdot1{,}1^3=133{,}1$ тыс. руб.

\major{3. Таблица: единый язык финансовой задачи}
Обозначим $a_i$ дополнительную операцию \emph{в начале} $i$-го года:
$a_i>0$ --- пополнение, $a_i<0$ --- снятие, $a_i=0$ --- операций нет.
Если на конец предыдущего года осталось $S_{i-1}$, то
\[
\underbrace{B_i}_{\substack{\text{начало года}\\\text{после операции}}}=S_{i-1}+a_i,\qquad
\underbrace{I_i}_{\text{проценты}}=B_ip_i,\qquad
\boxed{S_i=(S_{i-1}+a_i)(1+p_i)}.
\]
Здесь $p_i=P_i/100$ --- доля годовой ставки, которая может меняться.

\begin{center}
\renewcommand{\arraystretch}{1.35}
\small
\begin{tabular}{@{}c r r r r@{}}
\toprule
Год & Операция $a_i$ & Начальная $B_i$ & Проценты $I_i$ & Конечная $S_i$\\
\midrule
1 & $0$ & $S$ & $Sp$ & $Sk$\\
2 & $a$ & $Sk+a$ & $(Sk+a)p$ & $(Sk+a)k$\\
3 & $a$ & $Sk^2+ak+a$ & $(Sk^2+ak+a)p$ & $Sk^3+ak^2+ak$\\
\bottomrule
\end{tabular}
\end{center}
\smallskip
\textbf{Важно:} таблица выше описывает постоянную ставку и одинаковую операцию $a$ \emph{в начале второго и третьего года}. В первый год вклад не пополняют.

\major{4. Трёхлетний вклад с одинаковыми взносами}
Из третьей строки таблицы следует компактная формула:
\[
\boxed{\quad S_3=Sk^3+a(k^2+k).\quad}
\]
Если $a>0$, деньги добавляют; если $a<0$, снимают. При $a=0$ получаем формулу сложных процентов $S_3=Sk^3$.

\minor{Разобранный пример: ежегодное пополнение}
Вклад $120$ млн руб. открыт на $3$ года под $15\%$ годовых. В начале второго и третьего года добавляют по $10$ млн руб. Значит, $k=1{,}15$, $a=10$:
\[
\begin{aligned}
S_1&=120\cdot1{,}15=138,\\
S_2&=(138+10)\cdot1{,}15=170{,}2,\\
S_3&=(170{,}2+10)\cdot1{,}15=\boxed{207{,}23}\ \mone{млн руб.}
\end{aligned}
\]
То же в буквах: $S_3=120\cdot1{,}15^3+10\cdot1{,}15^2+10\cdot1{,}15=207{,}23$.

\minor{Разобранный пример: ежегодное снятие}
Вклад $100$ тыс. руб. открыт на $4$ года под $10\%$. С начала второго года ежегодно \emph{снимают} по $5$ тыс. руб.:
\[
\begin{aligned}
S_1&=100\cdot1{,}1=110,\\
S_2&=(110-5)\cdot1{,}1=115{,}5,\\
S_3&=(115{,}5-5)\cdot1{,}1=121{,}55,\\
S_4&=(121{,}55-5)\cdot1{,}1=\boxed{128{,}205}\ \mone{тыс. руб.}
\end{aligned}
\]
\textbf{Проверка смысла:} снятие учитывается \emph{сначала}, начисление процентов --- \emph{после}. Нельзя начислить проценты на уже снятые деньги.

\major{5. Как восстановить первоначальный вклад}
Три года, ставка $20\%$, в начале второго и третьего года вносят по $10$ тыс. руб. К концу третьего года на счёте $173{,}28$ тыс. руб. Пусть начальная сумма равна $x$ тыс. руб. Тогда
\[
1{,}2^3x+10\cdot1{,}2^2+10\cdot1{,}2=173{,}28,
\]
\[
1{,}728x+14{,}4+12=173{,}28,
\qquad x=\frac{146{,}88}{1{,}728}=\boxed{85}.
\]
\textbf{Смысл степеней:} исходная сумма приносит проценты три года; пополнение второго года --- два года; третьего --- один год.

\major{6. Как найти ставку и срок}
\minor{Неизвестна ставка}
Вложили $500$ тыс. руб.; через три года без пополнений получили $665{,}5$ тыс. руб.:
\[
500k^3=665{,}5,\qquad k^3=\frac{1331}{1000}=\left(\frac{11}{10}\right)^3.
\]
Отсюда $k=1{,}1$, а $P=100(k-1)=\boxed{10\%}$.

\minor{Неизвестен срок удвоения}
При ставке $10\%$ для удвоения суммы требуется $Sk^n\ge 2S$ ($S>0$):
\[
1{,}1^n\ge2,\qquad n\ge\frac{\lg 2}{\lg 1{,}1}\approx7{,}27.
\]
Проценты начисляются ежегодно, поэтому \textbf{минимальный целый срок --- $8$ лет}. Округляем вверх, поскольку через $7$ полных лет вклад ещё меньше удвоенного.

\major{7. Постоянная или переменная ставка?}
Если ставка не меняется и нет операций, достаточно $S_n=Sk^n$.
Если ставки по годам различаются, используйте таблицу или последовательное умножение:
\[
S_n=S\,(1+p_1)(1+p_2)\cdots(1+p_n)
\quad\mone{(при отсутствии дополнительных операций).}
\]
Например, при ставках $8\%$, $10\%$, $12\%$ за три года множитель равен
$1{,}08\cdot1{,}10\cdot1{,}12=1{,}33056$.
Для постоянных $10\%$ за три года множитель $1{,}1^3=1{,}331$.

\minor{Таблица для меняющейся ставки}
Пусть изначально внесли $1000$ тыс. руб., без дополнительных операций. Вариант со ставками $8\%$, $10\%$, $12\%$ оформим так:
\begin{center}
\renewcommand{\arraystretch}{1.3}
\begin{tabularx}{\linewidth}{@{}c X X X@{}}
\toprule
Год & В начале года & Множитель & В конце года \\
\midrule
1 & $1000$ & $1{,}08$ & $1080$ \\
2 & $1080$ & $1{,}10$ & $1188$ \\
3 & $1188$ & $1{,}12$ & $1330{,}56$ \\
\bottomrule
\end{tabularx}
\end{center}
Все суммы приведены в тысячах рублей. При ставке $10\%$ во все три года получилось бы $1331$ тыс. руб. --- на $440$ рублей больше.

\major{8. Частые ошибки и самопроверка}
\begin{itemize}
\item Ставка $15\%$ означает множитель $1{,}15$, а не $15$ и не $0{,}15$ для \emph{всей} конечной суммы.
\item В начале нового года операция меняет базу, на которую начисляется процент.
\item В трёхлетнем вкладе добавление в начале третьего года умножается на $k$, а не на $k^0$.
\item При переменной ставке нельзя подставить одну произвольную ставку в $Sk^n$.
\item При поиске числа \emph{полных} лет округляйте результат вверх, если требуется достижение заданной суммы.
\item При снятии проверьте допустимость: сумма после снятия должна быть неотрицательной.
\end{itemize}

\clearpage
\thispagestyle{fancy}
\begin{center}
{\Large\bfseries\color{darkaccent}Алгоритм: выбор способа решения}\par
\vspace{5mm}
{\small От условия задачи --- к наиболее короткой математической модели}
\end{center}
\vspace{7mm}
\begin{center}
\tikzset{
 alg/.style={draw=rulegray,rounded corners=2mm,very thick,align=center,text width=60mm,inner sep=6pt,font=\small,fill=white},
 dec/.style={diamond,aspect=2.55,draw=accent,very thick,align=center,text width=57mm,inner xsep=0pt,inner ysep=5pt,font=\small,fill=white},
 oval/.style={ellipse,draw=accent,very thick,align=center,text width=65mm,inner sep=5pt,font=\small,fill=white},
 arr/.style={-{Stealth[length=3mm,width=2mm]},thick,color=darkaccent}
}
\begin{tikzpicture}[node distance=14mm]
\node[oval] (begin) {Прочитайте условие.\\Определите момент изменения суммы};
\node[dec,below=17mm of begin] (choose) {Ставка постоянна\\и нет пополнений\\или снятий?};
\node[alg,below left=24mm and 1mm of choose,anchor=north] (simple) {\textbf{Да: готовая формула}\\[2mm]$S_n=S(1+P/100)^n$};
\node[alg,below right=24mm and 1mm of choose,anchor=north] (table) {\textbf{Нет: таблица по годам}\\[2mm]$B_i=S_{i-1}+a_i$\\$S_i=B_i(1+P_i/100)$};
\node[alg,below=26mm of $(simple.south)!0.5!(table.south)$] (calc) {Найдите неизвестную величину.\\Проверьте годы, знак операции, единицы измерения};
\node[oval,below=13mm of calc] (finish) {Запишите ответ};
\draw[arr] (begin)--(choose);
\draw[arr] (choose.west) -| node[pos=.24,above,font=\small]{да} (simple.north);
\draw[arr] (choose.east) -| node[pos=.24,above,font=\small]{нет} (table.north);
\draw[arr] (simple.south) |- (calc.west);
\draw[arr] (table.south) |- (calc.east);
\draw[arr] (calc)--(finish);
\end{tikzpicture}
\end{center}
\vspace{5mm}
\noindent\textbf{Полезная развилка.} Если ставка постоянна и пополнения одинаковы, для трёх лет можно сразу использовать $S_3=Sk^3+a(k^2+k)$; это частный случай расчёта по годам.

\clearpage
\major{Тренировочный блок --- Путь А}
{\color{textgray}10 задач, от базовых к более содержательным. В ответе указывайте единицы измерения. При необходимости расписывайте расчёт по годам.}
\vspace{3mm}
\begin{enumerate}[label=\textbf{\arabic*.},itemsep=9pt]
\item На вклад внесли $100$ тыс. руб. под $10\%$ годовых на $3$ года. Найдите сумму в конце срока при ежегодной капитализации и отсутствии дополнительных операций.
\item На вклад внесли $125$ тыс. руб. под $20\%$ годовых на $2$ года. Найдите конечную сумму.
\item На вклад внесли $200$ тыс. руб. под $5\%$ годовых на $3$ года. Найдите сумму на конец третьего года.
\item Вклад $120$ млн руб. открыт на $3$ года под $15\%$ годовых. В начале второго и третьего года вклад пополняют на $10$ млн руб. Найдите сумму в конце третьего года.
\item Вклад $100$ тыс. руб. открыт на $4$ года под $10\%$ годовых. В начале каждого года, начиная со второго, снимают по $5$ тыс. руб. Найдите конечную сумму.
\item На вклад сроком $3$ года под $20\%$ годовых внесли неизвестную сумму. В начале второго и третьего года добавили по $10$ тыс. руб. В конце третьего года на вкладе оказалось $173{,}28$ тыс. руб. Определите первоначальную сумму.
\item Через $3$ года вклад $500$ тыс. руб. без дополнительных операций вырос до $665{,}5$ тыс. руб. Найдите постоянную годовую процентную ставку.
\item Вклад ежегодно увеличивается на $10\%$ без пополнений и снятий. Через какое \emph{минимальное целое число лет} сумма впервые станет не меньше удвоенной первоначальной суммы?
\item Сравните два вклада по $1000$ тыс. руб. на $3$ года без дополнительных операций. Вариант А: ставка $10\%$ ежегодно. Вариант Б: ставки по годам $8\%$, $10\%$ и $12\%$. Найдите конечную сумму по каждому варианту и разницу между ними.
\item Вклад $200$ тыс. руб. размещён на $3$ года. Ставки по годам: $10\%$, $15\%$ и $5\%$. В начале второго года внесли ещё $10$ тыс. руб., а в начале третьего сняли $5$ тыс. руб. Найдите конечную сумму.
\end{enumerate}

\clearpage
\major{Ответы к тренировочному блоку}
\renewcommand{\arraystretch}{1.45}
\begin{center}
\begin{tabularx}{\linewidth}{@{}c X@{}}
\toprule
\textbf{№} & \textbf{Ответ}\\
\midrule
1 & $133{,}1$ тыс. руб.\\
2 & $180$ тыс. руб.\\
3 & $231{,}525$ тыс. руб.\\
4 & $207{,}23$ млн руб.\\
5 & $128{,}205$ тыс. руб.\\
6 & $85$ тыс. руб.\\
7 & $10\%$ годовых\\
8 & $8$ лет\\
9 & А: $1331$ тыс. руб.; Б: $1330{,}56$ тыс. руб.;\\
  & вариант А больше на $0{,}44$ тыс. руб. ($440$ руб.)\\
10 & $272{,}475$ тыс. руб.\\
\bottomrule
\end{tabularx}
\end{center}
\vspace{8mm}
\noindent{\color{darkaccent}\bfseries Контрольный вопрос}\par
\emph{Почему пополнение в начале второго года при трёхлетнем вкладе умножается на $k^2$, а пополнение в начале третьего --- на $k$?}
\vspace{5mm}
\noindent\color{textgray}Ответьте словами, используя число оставшихся начислений процентов.
\end{document}
